\section{Implicaciones Más Amplias del Estudio}

Los hallazgos de esta tesis trascienden el problema técnico concreto y tocan cuestiones de confianza, ética y regulación que conviene discutir de forma explícita. Las tres subsecciones siguientes exploran estas implicaciones más amplias, y muestran que la caracterización multidimensional de la calidad explicativa tiene consecuencias que van más allá de la elección de un método.

\subsection{Explicabilidad y Confianza en los Sistemas de Inteligencia Artificial}

La confianza que un analista o un regulador deposita en un sistema de detección no se sustenta únicamente en la exactitud de sus predicciones, sino también en la calidad de las razones que ofrece para justificarlas. Cuando un modelo alcanza un desempeño predictivo elevado pero acompaña cada alerta con explicaciones que cambian de una ejecución a otra, el usuario percibe que el sistema carece de un criterio consistente y termina desconfiando incluso de sus aciertos. Esta observación revela una distinción que resulta central para esta tesis: confiar en la predicción y confiar en la explicación son actos separados, gobernados por evidencias distintas. Un analista puede aceptar que una transacción es ilícita y, al mismo tiempo, rechazar la narrativa que el explicador construye alrededor de esa decisión, porque esa narrativa no le resulta ni estable ni verificable.

Los hallazgos de esta investigación refuerzan esa separación al mostrar que la estabilidad, la plausibilidad y la fidelidad constituyen dimensiones independientes de la calidad explicativa. Una explicación inestable erosiona la confianza porque impide que el analista construya un modelo mental duradero del comportamiento del sistema, ya que las mismas entradas producen justificaciones que fluctúan sin un patrón discernible. Una explicación infiel resulta aún más peligrosa, dado que ofrece al usuario una historia coherente y persuasiva que, sin embargo, no corresponde al razonamiento real del modelo. En ambos casos la precisión del clasificador queda subordinada a la credibilidad de su interfaz explicativa, porque el usuario humano actúa sobre lo que comprende y no sobre lo que el modelo calcula internamente.

Reconocer que estas tres propiedades no covarían tiene consecuencias directas sobre cómo se debe cultivar la confianza en entornos de auditoría financiera. Un explicador que produce mapas visualmente convincentes puede generar una confianza injustificada si esos mapas no reflejan las variables que efectivamente determinan la salida del modelo. Por el contrario, un explicador estable y fiel, aunque ofrezca explicaciones menos intuitivas, sienta una base más sólida para que el analista calibre cuándo delegar en el sistema y cuándo revisar manualmente. De este modo, la calidad de las explicaciones se convierte en un requisito para que la confianza sea legítima y no meramente aparente, y esta tesis argumenta que dicha calidad debe evaluarse de forma multidimensional en lugar de reducirse a una sola cifra agregada.

\subsection{Consideraciones Éticas de la Detección Automatizada}

Un sistema capaz de bloquear cuentas o de generar reportes con consecuencias legales opera sobre la vida de personas concretas, lo que traslada la discusión desde el plano técnico hacia el terreno de la responsabilidad ética. Cada alerta que el modelo produce puede desencadenar la congelación de fondos, la apertura de una investigación o la exclusión de un usuario del sistema financiero, medidas que afectan de manera desproporcionada a quienes resultan señalados por error. El costo humano de los falsos positivos no se distribuye de forma uniforme, puesto que las poblaciones con menor capacidad de reclamación suelen ser las que más difícilmente revierten una decisión automatizada. Por ello, el diseño de estos sistemas no puede tratar los falsos positivos como un simple parámetro a optimizar, sino como un daño real que exige mecanismos de reparación y revisión.

En este contexto, el derecho de una persona a recibir una explicación comprensible sobre por qué fue señalada adquiere un peso ético que trasciende la conveniencia operativa. Una explicación que el afectado y el auditor puedan interpretar es condición para que exista un debido proceso, dado que nadie puede impugnar una decisión cuyo fundamento permanece oculto o resulta ininteligible. La fidelidad de la explicación se vuelve entonces un asunto de justicia y no solo de calidad técnica, porque una justificación que no corresponde al razonamiento real del modelo priva al afectado de la posibilidad de defenderse frente a los motivos verdaderos de la decisión. Cuando la explicación es infiel, el proceso de apelación se construye sobre una ficción, y la garantía de defensa se transforma en una formalidad vacía.

A esta preocupación se suma el riesgo de que el modelo apoye sus decisiones en características espurias que correlacionan con la etiqueta sin guardar relación causal con la conducta ilícita. Si un clasificador aprende a asociar la ilicitud con atributos accidentales del conjunto de datos, puede penalizar sistemáticamente a usuarios que comparten esos atributos sin haber cometido ninguna irregularidad. La fidelidad explicativa cumple aquí una función de vigilancia, ya que solo una explicación que revele el razonamiento efectivo del modelo permite detectar cuándo este se apoya en señales indebidas. Por consiguiente, esta tesis sostiene que la evaluación de la fidelidad no es un ejercicio académico, sino una salvaguarda necesaria para que la automatización no reproduzca ni amplifique sesgos con impacto sobre personas inocentes.

\subsection{El Papel de la Explicabilidad en la Adopción Regulatoria de la IA}

La incorporación de modelos complejos en el sector financiero depende, en última instancia, de que los reguladores acepten decisiones cuyo mecanismo interno no siempre es transparente. Las autoridades de supervisión exigen que las entidades puedan justificar sus determinaciones ante auditorías, lo que coloca a la explicabilidad en el centro del proceso de adopción tecnológica. Existe, sin embargo, una tensión persistente entre el desempeño de los modelos más opacos y la exigencia regulatoria de transparencia, dado que las arquitecturas que capturan relaciones complejas suelen ser también las más difíciles de interpretar. Esta tensión explica por qué muchas instituciones prefieren modelos menos potentes pero legibles, sacrificando capacidad predictiva para preservar la posibilidad de rendir cuentas.

La contribución de esta tesis a ese debate consiste en mostrar que la calidad explicativa puede caracterizarse de forma estructurada a través de sus tres dimensiones, en lugar de tratarse como una propiedad monolítica que un método posee o no posee. Al distinguir la estabilidad, la plausibilidad y la fidelidad, el marco propuesto ofrece a los reguladores un vocabulario más preciso para especificar qué tipo de garantía explicativa requiere cada uso, ya que una auditoría orientada a la reparación de daños no demanda lo mismo que una orientada a la validación científica del modelo. Este enfoque también clarifica que ningún explicador optimiza simultáneamente las tres dimensiones, de modo que la selección del método debe subordinarse al propósito concreto de la supervisión. De esta manera, la elección técnica deja de ser arbitraria y pasa a responder a criterios que un regulador puede formular y verificar.

Caracterizar la calidad explicativa en estos términos ayuda a construir un puente entre la capacidad técnica de los modelos y su aceptación normativa, porque traduce una preocupación abstracta sobre la transparencia en requisitos evaluables y comparables. Cuando el regulador dispone de dimensiones separadas para exigir estabilidad, plausibilidad o fidelidad según el escenario, la conversación entre desarrolladores y supervisores puede apoyarse en criterios compartidos en lugar de en apreciaciones subjetivas. Esta claridad reduce la incertidumbre que frena la adopción, dado que una entidad sabe de antemano qué propiedad explicativa deberá demostrar para que un modelo complejo resulte admisible. En síntesis, esta tesis defiende que el avance hacia sistemas más potentes en el ámbito financiero no pasa por renunciar a la transparencia, sino por descomponerla en dimensiones que la técnica pueda entregar y que la regulación pueda auditar.
